% !TeX root = ../papers/nc-p3.tex
\section{被約な特異四次点曲線の被覆と多角形の係数判定}
\label{sec:reduced-quartic-point-curves}
\begin{theorem}\label{thm:reduced-quartic-point-curves}
四辺の左右が点ファイバーで、点グラフが純な被約四次曲線 $Y$ と交わらない有限部の和である
標準有限Frobenius窓を取る。滑らかな楕円四次のほかに残り得る $Y$ は、
二つの節点で交わる二つの平面二次曲線、または四辺四頂点の空間四角形である。
いずれも $(2,2)$ 完全交叉で、有限部は長さ4である。多角形では四移動は固定点のない回転を誘導し
（成分上の座標倍率を許す）、二階差は $K^2=1$ を満たす。
\end{theorem}
\begin{proof}
非連結なら全連結成分が次数3以下となる。合計が4なので次数1または2の成分を含む。
滑らかな開閉有理成分、または節点が一本しかない有理成分は固定微分の次数式で除かれる。
複数の成分を四移動が入れ替える場合にもこの式は成立する。

一位置で平面四次なら算術種数3である。一方、連結被約な非退化空間四次の算術種数は高々1である。
実際、$H^1(I_Y)=0$ なので一般平面切断の相異なる四点は平面を張る。
そのHilbert関数は次数 $0,1,2,\ldots$ で $1,3,4,\ldots$ である（四点の各々を他の三点から二直線で分離できる）。
非零因子となる平面式による斉次座標環の商はこの切断の座標環に全射するから、
四次曲線のHilbert多項式は $4n$ 以上となる。ゆえに一位置で平面なら全位置で平面であり、
隣接平面の積で一般の三つの双一次関係を取ると完全交叉曲線になる。共通基点が一次元以下だからである。
その各射影次数は3で四次曲線を含めない。

整で特異な非退化四次なら算術種数1、正規化は有理、全 $\delta$ は1となる。
正規化への余長1の局所環は二枝の共通値を課す節点か、一枝の $k+t^2k[[t]]$ の尖点である。
従ってGorensteinであり、双対化束は次数零の非零切断を持つので自明である。
前節の平面三次と同じ固定微分の証明で、唯一の特異点が固定されることに矛盾する。

平面三次部分曲線と外の直線の和も除ける。交点はただ一点で、平面方程式が直線上で単純に消える。
交点が三次の滑らかな点なら直線は節点一本の端となる。特異点なら局所座標で
平面を $z=0$、直線を $x=y=0$、三次の局所式を $f(x,y)$ として
和集合の式は $zx,zy,f$ である。直線上の余接層には長さ2の捩れがあり、
$B=0$、他の特異点はこの直線に乗らない。直線を点ごとに保つことも、
定理\ref{thm:fixed-differential}の第一項の零次の固定因子がこの交点を含むこともできない。
実際ここは全曲線で唯一接次元3の点なので固定される。

残りの成分次数を列挙する。
\begin{itemize}
\item $(3,1)$ の非平面三次は有理正規三次であり、直線との交わりは長さ2以下である
（この三次は二次式で切られる）。長さ1なら端の節点である。長さ2なら異なる二点での横断か一つの接点である。
二つの滑らかな有理成分を長さ2で貼った曲線の双対化束は自明である。
実際二点の場合は両側で留数が逆になる対数微分が生成し、一点の場合は通常の二枝接点となる。
後者は共通の一次jetを持つ二つの形式冪級数を貼った平面特異点で、導手の逆像は各枝上で次数2、
算術種数1である。各成分上で双対化束の次数は零で、非零の大域切断は
片方で恒等零なら接点で他方の定数も零となるため生成元となる。
四移動は以下このような自明な双対化束の有理成分の和では固定点を持てない
（定理\ref{thm:fixed-differential}、$\deg M<0$、Cartier固定因子の総次数零）。
ところが次数の異なる二成分を保つ自己同型は有理成分上で固定点を持つ。
\item $(2,2)$ では二平面の共通直線上に二曲線の交点があり、長さは高々2。
空、単純一交点、重複一交点は既述の議論で除ける。二単純交点なら記載した二角形である。
\item $(2,1,1)$ では三次部分曲線が平面なら既に除いた。従って各直線は二次曲線と高々一点で単純に交わる。
三つの対交点が相異なれば三角形の双対グラフで、双対化束は自明だが次数2の成分を保たねばならない。
対交点が不足すれば節点一本以下の端がある。三枝が同一点で会う場合は、どの二枝も接方向が異なる。
この三重点での各枝の余接捩れは長さ2である。三方向が独立なら三座標軸の和に形式同型である
（各枝をパラメータ表示し、その三つの表示のベクトル和による形式座標変更を用いる）。
三方向が一平面に入ればその平面への射影の三枝は通常の平面三重点で、
高さ座標はこの三枝上で二次以上から始まるため平面三重点の局所環に属する。
実際一枝にだけ二次以上の任意の級数を載せるには他の二枝の局所式の積を掛ければよい。
平面三枝の積の勾配は各枝でちょうど二次の共通因子を持つ。
いずれも次数式で恒等は除かれ、非恒等なら唯一の三重点が固定されることに反する。
\item 四直線では三本が同一平面に入る場合を既に除いた。四本共点なら方向は三点共線のない平面四点で、
二次式二本の完全交叉である。この錐の双対化束は自明、唯一の特異点が固定されてしまう。
三本だけが共点なら残る直線はその三本の高々一本としか会えず、端になってしまう。
三重点がなければ単純な交点グラフは三角形を含まない（相異なる三交点を作る三直線は同一平面に入る）。
全頂点の価数が2以上であることから四角形だけが残る。
\end{itemize}
多角形は節点の各枝に留数 $+1,-1$ の微分を順に貼って $\omega_Y=\mathcal O_Y$ を得る。
補題\ref{lem:genus-one-quartic-mark}から $(2,2)$ 完全交叉である。グラフの埋込みも局所完全交叉である
（片側で余次元2の完全交叉、他の三アフィン座標は第一側の正則関数の持上げで消去する）。
超過交叉公式で曲線の寄与は16となる。四移動に固定点はなく、向きを反す多角形自己同型は
成分か節点を固定するから、成分置換は非自明な回転となる。固定切断が消えないことから $K^2=1$ を得る。
\end{proof}

\begin{corollary}\label{cor:quartic-cycle-atlas}
二角形と四角形のそれぞれで、全係数に対するAS性、添字保持同型および最小周期は有限に判定できる。
$m=2,4$、$d=4/m$ とし、$m$ 本の射影直線の $\infty$ と次の $0$ を貼ってモデルを固定する。
多重次数 $(d,\ldots,d)$、貼合せ倍率1の束を $L$ とする。
$D\in k^*$、$K\in\{1,-1\}\subset\operatorname{Pic}^0(Y)=k^*$ に対し
$L_0=L,L_1=LD,L_2=LD^2K$ を取る。最初の二辺の切断乗法核（各次元8）から六平面を一つずつ選ぶ。
系\ref{cor:elliptic-point-finite-atlas}の二窓復元・点ファイバー検査に加え、第一グラフが
$Y$ と交わらない有限部の和となることを課す。これはグラフのイデアルを $I_Y$ で飽和して残余を取り、
その次元、非交叉とイデアルの交わりの一致を射影上で検査できる。
最後に定理\ref{thm:finite-orbit-quartic-marks}の四位置の二次元二次核の積を検査する。

同型比較は多角形の回転・反転と座標倍率 $z_a\mapsto s_a z_a$、
$(\prod_a s_a)^d=1$ の引戻しで二束・二関係を対応させることで必要十分である。
多角形の二種類の間に同型はない。移動後に最初の束を $L$ に戻して同じ比較を行えば
最小周期 $1,2,4$ を区別できる。
\end{corollary}
\begin{proof}
全切断系と乗法核の次元は補題\ref{lem:genus-one-quartic-mark}による。
各成分で束を標準自明化すると、貼合せの唯一の不変量は倍率の積である。
座標倍率を作用させると次数 $d$ の座標の両端の値の比が $s_a^d$ 倍となるので、
任意の第一束を $L$ に移せ、その安定群が記載のものになる。反転は $D$ を逆数にする。
曲線は点グラフの唯一の純曲線全体なので同型がこれを保ち、全切断系のため文字写像は引戻しで決まる。
逆は二窓の延長定理である。正則性は印の軌道が一点の有限判定を用いた。

連続な安定群の商も全ての零係数パターンで単項式比較に直せる。
連結トーラス $\prod s_a=1$ は各切断空間に対角的に作用する。
節点値を一つだけ立てる隣接二成分の端単項式と、各成分内部の単項式を基底に取ればよい。
貼合せを保つ各成分の倍率補正は、一成分で1と決め、次に $s_a^d$ を順に掛けて得る指標である。
両Grassmann座標を周囲の文字テンソル基底で書き、各非零Plücker支持で最初の座標を1とする。
残った非零座標の指標差の整数行列をSmith形にすれば、独立な指標の座標は全て1にでき、
残る単項式座標の一致が連結トーラス同型と同値である。代数閉体なので非零の整数乗写像は全射である。
最後に $\prod s_a\in\mu_d$ と多角形の有限置換を比較すればよい。
\end{proof}
